\cotefile{72}

\scanpage{98}

\uline{Dém.}

(1°) Pour la formule (5), qui se réduit à

\[s_i\in\unsure{E_i|\mathcal E_i}\qquad 0\leq i\leq n+1\]

on procède par réc. sur $i$, la formule étant triviale pour $i=0$. On prouve en $=$ temps (6), (qui est \unsure{vide|vides|vraie} pour $i=0$ si on pose $a_{-1}=s_0$, $\lambda_{-1}=1$, $E_{-1}=0$). On a

\[\begin{aligned}s_{i+1}&=\tau_i s_i=\underset{\substack{\rotatebox{-90}{$\in$}\\ E_i}}{s_i}+\langle s_i,a'_i\rangle a_i\in E_i+k a_i=E_{i+1},\\s_{i+1}&\equiv\langle s_i,a'_i\rangle a_i\quad(E_i)\end{aligned}\]

Comme

\[s_i\equiv\deleted{\unsure{\ell|k}}\lambda_{i-1}a_{i-1}\quad(E_{i-1})\]

on aura

\[\langle s_i,a'_i\rangle=\lambda_{i-1}\langle a_{i-1},a'_i\rangle=\lambda_i,\quad\deleted{\unsure{E}}\]

d'où la formule $(*)$ pour $i+1$ au lieu de $i$, si on suppose $a'_i\mid E_{i-1}=0$, OK.

\uline{Corollaire} Soient

\[\left\{\begin{aligned}h_i&=k a_i\\ H_i&=\operatorname{Ker}a'_i\end{aligned}\right.\quad\bigg|\quad\deleted{\unsure{0}}-1\leq i\leq n\deleted{\unsure{,\ h_{-1}=ks}.}\]

\margin{\uline{NB} $h_i,H_i$ \deleted{\unsure{ont|sont}} \unsure{définis|défini} en termes de $\tau_i$, \uline{dès que} $\{\unsure{a_i|s_i}\}$ est libre}

les sous-variétés associés à $\tau_i$, posons aussi $\left\{\begin{aligned}a_{-1}&=s_0\\h_{-1}&=ks_0\\[-2pt]&\deleted{\unsure{a_{-1}=s_0}}\end{aligned}\right.$. Les conditions \deleted{équivalentes} ci-dessous (A$'$) et (B$'$) sont équivalentes et impliquent la formule (15) $(E_i=k(s_0,\ldots,s_{i-1}))$

(A$'$) On a

\[\deleted{\unsure{(A')}}\quad h_i\subset H_j\quad\text{si}\quad -1\leq i<j\leq n,\quad j\ne i+1\tag*{(15$'$)}\]

\[\begin{aligned}&\deleted{\unsure{E_i=k(s_0,\ldots,s_{i})\quad 0\leq i\leq n+1}}\\ & (s_i)_{0\leq i\leq n+1}\text{ est libre }\bigl(\uline{\text{ou encore}},\ \deleted{\unsure{\text{équivalent}}}\\ &\deleted{\unsure{\text{il suffit de les}}}\ \deleted{\unsure{\Leftrightarrow H_j,\ h_j,\ i=n+1}}\ \deleted{\unsure{i,j}}\ h_i\not\subset \unsure{H'_{i+1}|H_{i+1}}\text{ en chaque}\\ &\text{point }-1\leq i\leq n-1\quad\text{indépendants}\bigr)\end{aligned}\tag*{(16$'$)}\]

(B$'$) On a

\[\langle a_i,a'_j\rangle=0\quad\text{pour}\quad-1\leq i<j\leq n,\quad j\ne i+1\tag*{(17)}\]

\[\deleted{\lambda_n\in k^*}\quad(a_i)_{-1\leq i\leq n}\text{ est libre}\tag*{(18)}\]

\[\lambda_n\in k^*\tag*{(19)}\]

\margin{\deleted{\unsure{$a_i$}}\quad\unsure{$0$|$a$}}

% margin note written sideways

En effet, (17) n'est qu'une reformulation de (15$'$), et il reste à prouver, en supposant (15$'$) et (17) vérifiées, que \unsure{(16$'$)|(15$'$)} équivaut à (18) + (19). \deleted{Pour $n=1$ \unsure{on c'est évident}.} On peut supposer que \deleted{\unsure{l'on a}} (si $n\geq1$) \inserted{$\lambda_{\unsure{n-1|n-2}}\ne0$}, c'est

% archive sheet number: 97
